\chapter{1912:Victor Grignard and Paul Sabatier}
\label{chap:chemistry1912}

The Nobel Prize in Chemistry for the year 1912 was awarded jointly to Victor Grignard and Paul Sabatier.

\textbf{Motivation:}
for the discovery of the so-called Grignard reagent, which in recent years has greatly advanced the progress of organic chemistry (Victor Grignard) and for his method of hydrogenating organic compounds in the presence of finely disintegrated metals whereby the progress of organic chemistry has been greatly advanced in recent years (Paul Sabatier)

\section{Victor Grignard}

\subsection*{Early Life and Education}

Victor Grignard was born on 1871-05-06 in Cherbourg, France.
François Auguste Victor Grignard was a French chemist who won the Nobel Prize for his discovery of the eponymously named Grignard reagent and Grignard reaction, both of which are important in the formation of carbon–carbon bonds. He also wrote some of his experiments in his laboratory notebooks.

\subsection*{Career and Major Research}

Grignard was professor at Nancy and later at Lyon. In 1900 he developed the organomagnesium compounds that bear his name, obtained by reacting magnesium with organic halides, and showed that these reagents combine readily with aldehydes and ketones to form alcohols. The Grignard reaction became one of the most versatile tools of organic synthesis.

\subsection*{Nobel Prize-winning Work}

The work for which Victor Grignard was awarded the Nobel Prize in Chemistry in 1912 was described by the Nobel Committee as: "for the discovery of the so-called Grignard reagent, which in recent years has greatly advanced the progress of organic chemistry".

This research fundamentally advanced the field by making groundbreaking contributions that transformed chemical science.

\subsection*{Significance and Impact}

The impact of Victor Grignard's work extends far beyond the specific achievement recognized by the Nobel Prize.
These contributions continue to influence chemical research and education today.

\subsection*{Later Career and Honors}

Victor Grignard continued his scientific work until 1935-12-13, passing away in Lyon, France.
He received numerous honors including membership in major scientific academies and societies.

\subsection*{Legacy}

The legacy of Victor Grignard endures through the lasting impact of his scientific contributions.
Victor Grignard's work continues to be cited and built upon by researchers across the chemical sciences.

\subsection*{Modern Developments}

Grignard reagents remain among the most widely used organometallic reagents in organic synthesis, employed in the manufacture of pharmaceuticals, agrochemicals, flavors, and polymers. The Grignard reaction is a staple of both academic laboratories and industrial process chemistry, and modern cross-coupling chemistry descends directly from this carbon-carbon bond-forming tool.

\subsection*{Selected Publications}

Victor Grignard's major publications include numerous papers in leading journals of the era, detailing the experimental and theoretical work summarized above.

\clearpage

\section{Paul Sabatier}

\subsection*{Early Life and Education}

Paul Sabatier was born on 1854-11-05 in Carcassonne, France.

\subsection*{Career and Major Research}

Sabatier was professor of chemistry at Toulouse. He showed that finely divided metals, especially nickel, catalyze the hydrogenation of unsaturated organic compounds, a method he applied to the conversion of oils into fats and to the hydrogenation of benzene. His work laid the foundation of industrial catalytic hydrogenation.

\subsection*{Nobel Prize-winning Work}

The work for which Paul Sabatier was awarded the Nobel Prize in Chemistry in 1912 was described by the Nobel Committee as: "for his method of hydrogenating organic compounds in the presence of finely disintegrated metals whereby the progress of organic chemistry has been greatly advanced in recent years".

This research fundamentally advanced the field by making groundbreaking contributions that transformed chemical science.

\subsection*{Significance and Impact}

The impact of Paul Sabatier's work extends far beyond the specific achievement recognized by the Nobel Prize.
These contributions continue to influence chemical research and education today.

\subsection*{Later Career and Honors}

Paul Sabatier continued his scientific work until 1941-08-14, passing away in Toulouse, France.
He received numerous honors including membership in major scientific academies and societies.

\subsection*{Legacy}

The legacy of Paul Sabatier endures through the lasting impact of his scientific contributions.
Paul Sabatier's work continues to be cited and built upon by researchers across the chemical sciences.

\subsection*{Modern Developments}

Sabatier's hydrogenation catalysis, developed in the early 1900s, remains fundamental to the modern hydrogenation industry, from margarine production and edible oil hardening to ammonia synthesis and the hydrogenation of fats, oils, and petrochemicals. His nickel-based catalysts inspired the heterogeneous catalysis at the heart of modern ammonia, methanol, and Fischer–Tropsch plants.

\subsection*{Selected Publications}

Paul Sabatier's major publications include numerous papers in leading journals of the era, detailing the experimental and theoretical work summarized above.

\clearpage

\section*{Chapter Summary}

Victor Grignard and Paul Sabatier shared the prize, Grignard for his organomagnesium reagents and Sabatier for his method of hydrogenating organic compounds with finely divided metals.
